\documentclass[10pt,a4paper]{amsart}
\usepackage[margin=19mm]{geometry}
\usepackage{amsmath,amssymb,mathtools}
\usepackage[scr=rsfs,cal=euler]{mathalfa}
\usepackage{xfrac}
\usepackage[unicode,hidelinks,bookmarks=false]{hyperref}
\newcommand{\ie}{i.e.\ }
\newcommand{\cf}{cf.\ }
\newcommand{\resp}{respectively\ }
\newcommand{\Subsection}{\subsection}
\providecommand{\nobreakdash}{\nobreak\hbox{-}}
\setlength{\emergencystretch}{2em}
\multlinegap0pt
\usepackage{bm}
\usepackage{bbm}
\usepackage{dsfont}
\usepackage{xspace}
\usepackage{cancel}
\usepackage{mathtools}
\usepackage{amsthm}
\makeatletter
\@ifundefined{theorem}{\newtheorem{theorem}{Theorem}}{}
\@ifundefined{lem}{\newtheorem{lem}[theorem]{Lemma}}{}
\makeatother
\newcommand{\ssubset}{\subset\joinrel\subset}
\title[The Levi $q$-core and Property ($P\_q$)]{The Levi $q$-core and Property ($P\_q$)}
\author{Gian Maria Dall'Ara, Samuele Mongodi, John N. Treuer}
\date{Source: arXiv:2505.17693v1 (CC BY 4.0)}
\begin{document}
\maketitle

\noindent\emph{Source and licence.} The Levi $q$-core and Property ($P_q$). Version arXiv:2505.17693v1 (CC BY 4.0), distributed under the Creative Commons Attribution 4.0 International licence, as recorded in the arXiv licence metadata for this exact version and shown on the abstract page https://arxiv.org/abs/2505.17693v1. Full rights notice: licenses/arxiv-2505.17693-CC-BY-4.0.txt.\par\smallskip

\noindent\textit{Editorial scope.} Counted unit: the Lem whose source label is \texttt{break up of the longer lemma into a second part} in the article (source0.tex lines 521--523, source label \texttt{break up of the longer lemma into a second part}), delivered together with the article's own complete original proof of it (source0.tex lines 524--565, one \texttt{proof} environment of 42 lines). No mathematics is written, repaired or re-derived: statement and proof are the article's own text line for line. Required context retained uncounted: closed under maximums (lines 307--309); gluing lemma (lines 318--328); distance function is Pq (lines 507--514). Every definition, display and label of those lines is retained unchanged. Omitted from this package and not redistributed: 1--306, 310--317, 329--506, 515--520, 566--668. Nothing inside a retained excerpt is omitted or altered: every retained body is delivered exactly as the article's lines are written, so no omission receipt of any kind accompanies this package. Citations: the retained excerpts carry no citation command at all, so no bibliography entry of the article is transcribed and no reference is redistributed. Reference adaptation: none was needed and none was made; every reference inside the delivered unit resolves inside it, so no unresolved reference or citation is produced. Printed slips kept literal: no printed word, symbol, hypothesis or number of the counted statement or of its proof was altered. Layout and class: the article's own class is \texttt{amsart} with packages including amsmath, amsfonts, amsthm, amssymb, color, hyperref, enumerate, extarrows, cite, graphicx, fontenc, backref; the wrapper is the lane's pinned \texttt{amsart} editorial wrapper at 10pt on A4, which restates the theorem-like environments the retained text needs and adds the article's own macro definitions that the retained text needs (\texttt{\textbackslash ssubset}), with extra wrapper packages (bm, bbm, dsfont, xspace, cancel, mathtools). The delivered numbering is the unit's own. Assets: no retained span contains a figure environment, a graphics-inclusion command, a PDF-page inclusion, a table or a TikZ picture; no image file is copied into this package, the package declares no asset dependency and no image file is redistributed with it. Licence: version arXiv:2505.17693v1 of this article is distributed under the Creative Commons Attribution 4.0 International licence, as recorded in the arXiv licence metadata for this exact version and shown on the abstract page https://arxiv.org/abs/2505.17693v1. A full rights notice is delivered at licenses/arxiv-2505.17693-CC-BY-4.0.txt. Characters: every retained excerpt is pure ASCII, so the delivered engine, XeLaTeX with its default Latin Modern text font, reports no missing character.
\begin{lem}\label{closed under maximums}
    If $U$ is an open set and $\phi_1, \phi_2 \in P_q(U)$, then $\max\{\phi_1, \phi_2\} \in P_q(U)$. 
\end{lem}

\begin{lem}[Gluing Lemma]\label{gluing lemma}
Let $U$ be an open set and let $\omega$ be a non-empty proper open subset of $U$.  If $u \in P_q(U)$, $v \in P_q(\omega)$, and $max\{u, v\} = u$ in a neighborhood in $U$ of each $y \in \partial \omega \cap U$, then the formula
$$
w = \begin{cases}
    \max\{u, v\} & \omega
    \\
    u & U\setminus\omega
\end{cases}
$$
defines a function in $P_q(U)$.
 \end{lem}

\begin{lem}\label{distance function is Pq}
Let $X$ be compact, $F = X \setminus \hbox{int }J_q(X)$ be nonempty and $h(z) = d(z, F)$.  Then $h \in P_q(X)$  Moreover, for any $\epsilon > 0$, there is a neighborhood $W$ of $X$ and a function $\psi \in P_q(W)$ such that 
$
|\psi(z) - h(z)| < \epsilon,
$
for all $z \in W$.

\end{lem}

\begin{lem}\label{break up of the longer lemma into a second part}
    Let $F = X \setminus \hbox{int }J_q(X)$.  If $U$ is a neighborhood of $F$ and $\phi \in P_q(U)$ with $\phi > 0$, then there is a function $\theta$ defined in a neighborhood $W$ of $X$ such that $\theta \in P_q(W)$ and $\theta = \phi$ on a neighborhood of $F$.
\end{lem}

\begin{proof}
After shrinking $U$ we may suppose that $\phi$ is bounded on $U$.  Let $\delta$ be such that 
\begin{equation}\label{how big U is}
\{d(z, F) < 4\delta\} \ssubset U.
\end{equation}
Using Lemma \ref{distance function is Pq}, let $W$ be a neighborhood of $X$ and $\psi_1$ be a function in $P_q(W)$ such that
$|\psi_1(z) - d(z, F)| < \delta$ for all $z \in W$. The set $W \cap \{d(z, F) < \delta\}$ is a nonempty neighborhood of $F$, and $W \cap \{d(z, F) \geq 3.5\delta\}$ is a possibly empty set.
% By intersecting $W$ with $\{d(z, X) < \delta\}$ we may suppose that
% $$
% W \subset \{d(z, X) < \delta\}.
% $$
% We note for later that for all $z$, $d(z, X) \leq d(z, F)$. 
Let $\psi_2:W \to \mathbb{R}$ by $\psi_2 = \psi_1 - 2\delta$.  Then
\begin{eqnarray*}
    \psi_2 = \psi_1 - 2\delta &<& d(z, F) + \delta - 2\delta < 0, \quad z \in W \cap \{d(z, F) < \delta\}.
\end{eqnarray*}
Additionally,
\begin{eqnarray*}
\psi_2(z) = \psi_1(z) - 2\delta &\geq& d(z, F) - \delta - 2\delta \geq .5\delta, \quad W \cap \{d(z, F) \geq 3.5\delta\}.
\end{eqnarray*}
Rescale $\psi_2$ such that
\begin{equation}\label{where psi2 is negative and positive}
\psi_2(z) < 0, \quad z \in \{d(z, F) < \delta\} \cap W, \quad \psi_2(z) > \max_{U}\phi, \quad W \cap \{d(z, F) \geq 3.5\delta\}.
\end{equation}
By the continuity of $\psi_2$, $\psi_2(z) > \max_U \phi$ holds in a neighborhood in $W$ of $W \cap \{d(z, F) \geq 3.5\delta\}$. Define $\theta$ by
\begin{equation}
\theta = \begin{cases} 
\max(\psi_2, \phi), & W\cap \{d(z, F) < 3.5\delta\}
\\
\psi_2 & W \cap \{d(z, F) \geq 3.5\delta\}.
\end{cases}
\end{equation}
By \eqref{how big U is}, $\phi$ is well-defined on the nonempty neighborhood $W\cap \{d(z, F) < 3.5\delta\}$.  If $W \cap \{d(z, F) \geq 3.5\delta\}$ is empty, then using Lemma \ref{closed under maximums}, $\theta \in P_q(W)$.  If it is nonempty, then by Lemma \ref{gluing lemma}, $\theta \in P_q(W)$.  
Since $\phi > 0$, by \eqref{where psi2 is negative and positive},
$
\theta|_{W \cap \{d(z, F) < \delta\}} = \phi.
$
 The proof of the lemma is complete.


    
\end{proof}

\end{document}
